\subsection{Proof of Theorem \ref{thm:invariant_measure}}\label{app:thm_invariant_measure}
Set \(\vartheta_0=\vartheta_\star\) and \(Z_0\sim\nu_{\theta_\star}\). We construct a solution to \eqref{eq:general_mf_system_opt}--\eqref{eq:general_mf_system_x} with initial law \(\delta_{\vartheta_\star}\otimes\nu_{\theta_\star}\) and show that this law is preserved.

Set \(\vartheta_t\equiv\vartheta_\star\). Since \(\Pi_\theta(\vartheta_\star)=\theta_\star\), let \(Z_t\) solve the frozen sampler equation
\[
    \mathrm dZ_t
    =
    b_{\theta_\star}(Z_t)\,\mathrm dt
    +
    \sigma_{\theta_\star}(Z_t)\,\mathrm dB_t,
    \qquad
    Z_0\sim\nu_{\theta_\star}.
\]
By invariance of \(\nu_{\theta_\star}\), \(\mathrm{Law}(Z_t)=\nu_{\theta_\star}\) for every \(t\ge0\). Hence
\begin{align*}
G_{\nu_t}(\vartheta_t)
&=
\int_{\sX\times\sV}
F(\Pi_\theta(\vartheta_t),\Pi_x(z))\,\nu_t(\mathrm dz)\\
&=
\int_{\sX\times\sV}
F(\theta_\star,\Pi_x(z))\,\nu_{\theta_\star}(\mathrm dz)\\
&=
\int_{\sX}F(\theta_\star,x)\,\pi_{\theta_\star}(\mathrm dx)
=
\nabla_\theta\ell(\theta_\star)
=0,
\end{align*}
where we used \((\Pi_x)_\#\nu_{\theta_\star}=\pi_{\theta_\star}\). Assumption~\Cref{ass:existence_minimizer} gives \(\Psi_t(\vartheta_\star,0)=0\), and therefore \(\mathrm d\vartheta_t=0\). Thus \((\vartheta_t,Z_t)\) solves the coupled system and has law \(\delta_{\vartheta_\star}\otimes\nu_{\theta_\star}\) for all \(t\ge0\). By pathwise uniqueness in Assumption~\ref{ass:strong_sol_exist_unique}, this constructed solution is the unique solution with the prescribed initial data. This proves that \(\mu_\star\) is invariant.

\subsection{Proof of Theorem \ref{thm:twisted_joint_strong_monotonicity}}\label{app:thm_joint_strong_monotonicity}

The proof follows a synchronous coupling argument with twisted metrics $D_\vartheta,D_z$, similar in style to Proposition 8 in \citet{theorem2literature_bolley_guillin_malrieu}.

Let \((\vartheta_t,\nu_t)\) and \((\tilde\vartheta_t,\tilde\nu_t)\) be two solutions. Choose an optimal coupling \(\Lambda_0\in\Gamma(\nu_0,\tilde\nu_0)\), take \((Z_0,\tilde Z_0)\sim\Lambda_0\), and drive the two sampler equations with the same Brownian motion:
\begin{align*}
    \mathrm dZ_t
    &=
    b_{\Pi_\theta(\vartheta_t)}(Z_t)\,\mathrm dt
    +
    \sigma_{\Pi_\theta(\vartheta_t)}(Z_t)\,\mathrm dB_t,\\
    \mathrm d\tilde Z_t
    &=
    b_{\Pi_\theta(\tilde\vartheta_t)}(\tilde Z_t)\,\mathrm dt
    +
    \sigma_{\Pi_\theta(\tilde\vartheta_t)}(\tilde Z_t)\,\mathrm dB_t.
\end{align*}
By uniqueness, the marginal laws of \(Z_t\) and \(\tilde Z_t\) are \(\nu_t\) and \(\tilde\nu_t\). Hence \(\Lambda_t:=\mathrm{Law}(Z_t,\tilde Z_t)\) is a coupling of \(\nu_t\) and \(\tilde\nu_t\).

Define
\[
    D_t
    :=
    \|\vartheta_t-\tilde\vartheta_t\|_{D_\vartheta}^2
    +
    \mathbb E\|Z_t-\tilde Z_t\|_{D_z}^2
\]
and let
\(\lambda_\vartheta^-,\lambda_\vartheta^+\) be the smallest and largest
eigenvalues of \(D_\vartheta\), and let \(\lambda_z^-,\lambda_z^+\) be the
corresponding eigenvalues of \({D_z}\). For the optimiser states,
\begin{align*}
    \frac{\mathrm d}{\mathrm dt}
    \|\vartheta_t-\tilde\vartheta_t\|_{D_\vartheta}^2
    =
    2\left\langle \vartheta_t-\tilde\vartheta_t,\,D_{\vartheta}\big(
    \Psi_t(\vartheta_t,G_{\nu_t}(\vartheta_t))
    -
    \Psi_t(\tilde\vartheta_t,G_{\tilde\nu_t}(\tilde\vartheta_t))\big)
    \right\rangle .
\end{align*}
For the sampler variables, let $Q_t:=\|Z_t-\tilde Z_t\|^2_{D_z}$ and introduce the notation,
\begin{equation*}
    \Delta Z_t=Z_t-\tilde Z_t,\qquad\Delta b_t=b_{\Pi_\theta(\vartheta_t)}(Z_t)
    -
    b_{\Pi_\theta(\tilde\vartheta_t)}(\tilde Z_t),\qquad \Delta\sigma_t=\sigma_{\Pi_\theta(\vartheta_t)}(Z_t)
    -
    \sigma_{\Pi_\theta(\tilde\vartheta_t)}(\tilde Z_t).
\end{equation*}
Then It\"o's formula under the synchronous coupling gives
\begin{align*}
    Q_t=Q_0+\int_0^t2\left\langle \Delta Z_s,D_z\Delta b_s
    \right\rangle + \| D^{1/2}_z\Delta \sigma_s \|^2_{\mathrm{HS}}\md s +M_t,\qquad M_t:= 2\int_0^t\left\langle D_z\Delta Z_s,\Delta \sigma_s\md B_s\right\rangle.
\end{align*}
If we assume under \ref{ass:strong_sol_exist_unique} that the diffusion coefficient is constant, $M_t=0$ and we can proceed to take expectations. Otherwise for any $t>0$ we pick a $T>t$ and show $M$ is a true martingale on $[0,T]$ and hence has zero expectation under the linear growth assumption. We proceed by bounding the expected square root of the quadratic variation under the linear-growth assumption in \ref{ass:strong_sol_exist_unique}. Indeed, the quadratic variation at time $T$ is given
\begin{align*}
    \langle M\rangle_T&=4\int_0^T\left\|(\Delta\sigma_s)^TD_z\Delta Z_s\right\|^2\md s \\
    &\le 4\int_0^T \left\|D^{1/2}_z\Delta\sigma_s\right\|^2_{\mathrm{HS}}\underbrace{\left\|D^{1/2}\Delta Z_s\right\|^2}_{=Q_s}\md s \\
    &\le 4T\left(\sup_{s\le T}\left\|D^{1/2}_z\Delta\sigma_s\right\|^2_{\mathrm{HS}}\right)\left(\sup_{s\le T} Q_s\right)
\end{align*}
hence taking expected square roots and applying Cauchy-Schwarz,
\begin{align*}
    \mathbb{E}\sqrt{\langle M\rangle_T}&\le 2\sqrt{T}\,\mathbb{E}\left[\left(\sup_{s\le T}\left\|D^{1/2}_z\Delta\sigma_s\right\|^2_{\mathrm{HS}}\right)^{1/2}\left(\sup_{s\le T} Q_s\right)^{1/2}\right]\\
    &\le 2\sqrt{T}\sqrt{\mathbb{E}\left[\sup_{s\le T}\left\|D^{1/2}_z\Delta\sigma_s\right\|^2_{\mathrm{HS}}\right]}\sqrt{\mathbb{E}\left[\sup_{s\le T} Q_s\right]}.
\end{align*}
Now for finite $T$, we may define $r_T:=\sup_{t\le T}\max\{\|\theta_t\|,\|\tilde\theta_t\|\}<\infty$ since $\theta_t,\tilde\theta_t$ are deterministic, continuous and bounded on $[0,T]$. Then on $[0,T]$ we have the modified linear-growth assumption,
\begin{equation}\label{eq:finite_horizon_linear_growth}
    \|b_{\Pi_{\theta}(\vartheta_s)}(Z_s)\|+\|\sigma_{\Pi_{\theta}(\vartheta_s)}(Z_s)\|\le K_T(1+\|Z_s\|),\qquad K_T:=K(1+r_T),
\end{equation}
which we may apply to bound
\begin{equation*}
    \left\|D^{1/2}_z\Delta\sigma_s\right\|^2_{\mathrm{HS}}\le 2\lambda^+_z\left(\|\sigma_{\Pi_\theta(\vartheta_s)}(Z_s)\|^2_{HS}+\|\sigma_{\Pi_\theta(\tilde\vartheta_s)}(\tilde Z_s)\|^2_{HS}\right)\le 4\lambda^+_z K_T^2\left(2+\|Z_s\|^2+\|\tilde Z_s\|^2\right),
\end{equation*}
and we also have
\begin{equation*}
    Q_s\le \lambda^+_z\|Z_s-\tilde Z_s\|^2\le 2\lambda_z^+\left(\|Z_s\|^2+\|\tilde Z_s\|^2\right).
\end{equation*}
The modified linear-growth assumption \eqref{eq:finite_horizon_linear_growth} also implies the expectations of the suprema of $\|Z_s\|^2$ and $\|\tilde Z_s\|^2$ are bounded (see for example Theorem 14.5.2 in \citet{Pascucci2024}). This gives the desired bound $\mathbb{E}\sqrt{\langle M\rangle_T}<\infty$ and so, by the Burkholder–Davis–Gundy inequality, $M_t$ is a true martingale on $[0,T]$. We can now take expectations then differentiate with respect to $t$ and obtain
\begin{align*}
    \frac{\mathrm d}{\mathrm dt}
    \mathbb E\|\Delta Z_t\|_{D_z}^2
    &=
    \mathbb E\Big[
    2\left\langle \Delta Z_t,\,
   D_z\Delta b_t
    \right\rangle+
    \left\|D_z^{1/2}\Delta\sigma_t\right\|_{\mathrm{HS}}^2
    \Big].
\end{align*}
Adding the two identities and applying \eqref{eq:joint_strong_monotonicity_twisted} with \(\Lambda=\Lambda_t\) yields
\[
    \frac{\mathrm d}{\mathrm dt}D_t
    \le
    -2\rho D_t\qquad\Rightarrow\qquad D_t\le e^{-2\rho t}D_0
\]
by Gronwall's lemma. Since \(\Lambda_0\) is optimal,
\begin{align*}
    D_0
    &=
    \|\vartheta_0-\tilde\vartheta_0\|_{D_\vartheta}^2
    + \mathbb E\|Z_0-\tilde Z_0\|_{D_z}^2\\
    &\le \max(\lambda_\vartheta^+,\lambda_z^+)\big(\|\vartheta_0-\tilde\vartheta_0\|^2
    + \mathbb E\|Z_0-\tilde Z_0\|^2\big)
    \\
    &= \max(\lambda_\vartheta^+,\lambda_z^+)\big(\|\vartheta_0-\tilde\vartheta_0\|^2
    +W_2^2(\nu_0,\tilde\nu_0)\big),
\end{align*}
and for every \(t\ge0\),
\begin{align*}
    \|\vartheta_t-\tilde\vartheta_t\|^2 +W_2^2(\nu_t,\tilde\nu_t)
    &\le
    \|\vartheta_t-\tilde\vartheta_t\|^2+\mathbb E\|Z_t-\tilde Z_t\|^2\\
    &\le\frac{1}{\min(\lambda^-_\vartheta,\lambda^-_z)}\big(\|\vartheta_t-\tilde\vartheta_t\|^2_{D_\vartheta}+\mathbb E\|Z_t-\tilde Z_t\|_{D_z}^2\big)
\end{align*}
we arrive at,
\begin{align*}
    \mathbf{d}((\vartheta_t,\nu_t),(\tilde{\vartheta}_t,\tilde{\nu}_t))^2
    \le
    \frac{\max(\lambda_\vartheta^+,\lambda_z^+)}{\min(\lambda^-_\vartheta,\lambda^-_z)}e^{-2\rho t}
    \mathbf{d}((\vartheta_0,\nu_0),(\tilde{\vartheta}_0,\tilde{\nu}_0))^2,
\end{align*}
proving the contraction result.

Now assume in addition that \ref{ass:existence_minimizer} holds.
By Theorem~\ref{thm:invariant_measure},
\((\vartheta_\star,\nu_{\theta_\star})\) is a stationary solution. Moreover,
for deterministic optimiser states,
\[
\begin{aligned}
W_2^2\!\left(
\delta_\vartheta\otimes\nu,\,
\delta_{\tilde\vartheta}\otimes\tilde\nu
\right)
&=
\|\vartheta-\tilde\vartheta\|^2
+
W_2^2(\nu,\tilde\nu).
\end{aligned}
\]
Applying the preceding contraction with
\((\tilde\vartheta_t,\tilde\nu_t)
=(\vartheta_\star,\nu_{\theta_\star})\) proves
\[
W_2(\mu_t,\mu_\star)
\leq
\sqrt\kappa\,e^{-\rho t}W_2(\mu_0,\mu_\star).
\]
Finally, let \(\delta_\vartheta\otimes\nu\) be any other stationary law of
the stated product form with finite second moment. Applying this estimate
between that stationary law and \(\mu_\star\) gives, for every \(t\geq0\),
\[
W_2(\delta_\vartheta\otimes\nu,\mu_\star)
\leq
\sqrt\kappa\,e^{-\rho t}
W_2(\delta_\vartheta\otimes\nu,\mu_\star).
\]
Letting \(t\to\infty\) shows that the left-hand side is zero. Hence
\(\delta_\vartheta\otimes\nu=\mu_\star\), proving uniqueness within the class asserted in the theorem.

\subsection{Proof of Theorem \ref{thm:finite_particle_error}}\label{app:proof_finite_particle_error}
We follow the synchronous coupling route using an $N$-copy mean-field system as in \citet{pgd_theory}. Introducing the notation
\[
\begin{aligned}
    \Delta\vartheta_t&:=\vartheta_t^N-\bar\vartheta_t,
    &
    \Delta Z_t^i&:=Z_t^i-\bar Z_t^i,\\
    \xi_t^N
    &:=
    \|\Delta\vartheta_t\|^2
    +
    \frac1N\sum_{i=1}^N\|\Delta Z_t^i\|^2,
    &
    \widehat\xi_t^N
    &:=
    \|\Delta\vartheta_t\|_{D_\vartheta}^2
    +
    \frac1N\sum_{i=1}^N\|\Delta Z_t^i\|_{{D_z}}^2.
\end{aligned}
\]
We also write $\theta^N_t=\Pi_\theta(\vartheta^N_t)$ and $\bar\theta_t=\Pi_\theta(\bar\vartheta_t)$. By It\"o's formula and the synchronous coupling,
\begin{align*}
    \md \widehat\xi_t^N
&= 
2\left\langle
    \Delta\vartheta_t,
    D_\vartheta
    \left[
    \Psi_t(\vartheta_t^N,G_{\nu_t^N}(\vartheta_t^N))
    -
    \Psi_t(\bar\vartheta_t,G_{\bar\nu_t}(\bar\vartheta_t))
    \right]
\right\rangle \md t\\
&\qquad+
\frac1N\sum_{i=1}^N
\Big\{
2\left\langle
    \Delta Z_t^i,
    {D_z}
    \left[
    b_{\theta^N_t}(Z_t^i)
    -
    b_{\bar\theta_t}(\bar Z_t^i)
    \right]
\right\rangle\\
&\qquad\qquad\qquad+
\left\|
    {D_z}^{1/2}
    \left[
    \sigma_{\theta^N_t}(Z_t^i)
    -
    \sigma_{\bar\theta_t}(\bar Z_t^i)
    \right]
\right\|_{\mathrm{HS}}^2
\Big\} \md t\\
&\qquad + \md M^N_t,\\
\md M^N_t&:=\frac{2}{N}\sum_{i=1}^N\left\langle D_z\Delta Z^i_t,\left[
    \sigma_{\theta^N_t}(Z_t^i)
    -
    \sigma_{\bar\theta_t}(\bar Z_t^i)
    \right]\md B^i_t\right\rangle,\qquad M^N_0=0.
\end{align*}

% \begin{align*}
% \frac{\mathrm d}{\mathrm dt}\mathbb E[\widehat\xi_t^N]
% &=
% \mathbb E\Bigg[
% 2\left\langle
%     \Delta\vartheta_t,
%     D_\vartheta
%     \left[
%     \Psi_t(\vartheta_t^N,G_{\nu_t^N}(\vartheta_t^N))
%     -
%     \Psi_t(\bar\vartheta_t,G_{\bar\nu_t}(\bar\vartheta_t))
%     \right]
% \right\rangle\\
% &\quad+
% \frac1N\sum_{i=1}^N
% \Big\{
% 2\left\langle
%     \Delta Z_t^i,
%     {D_z}
%     \left[
%     b_{\Pi_\theta(\vartheta_t^N)}(Z_t^i)
%     -
%     b_{\Pi_\theta(\bar\vartheta_t)}(\bar Z_t^i)
%     \right]
% \right\rangle\\
% &\qquad\qquad+
% \left\|
%     {D_z}^{1/2}
%     \left[
%     \sigma_{\Pi_\theta(\vartheta_t^N)}(Z_t^i)
%     -
%     \sigma_{\Pi_\theta(\bar\vartheta_t)}(\bar Z_t^i)
%     \right]
% \right\|_{\mathrm{HS}}^2
% \Big\}
% \Bigg].
% \end{align*}
Note that since solutions exist, the process $M^N_t$ is a local martingale. Now add and subtract
\(D_\vartheta\Psi_t(\bar\vartheta_t,
G_{\bar\nu_t^{[N]}}(\bar\vartheta_t))\). Since
Assumption~\Cref{ass:twisted_joint_strong_monotonicity} holds for every choice of states, measures and couplings, we may apply it to the empirical
coupling
\[
    \Lambda_t^N
    :=
    \frac1N\sum_{i=1}^N
    \delta_{(Z_t^i,\bar Z_t^i)}
\]
giving
\begin{align*}
    \md \hat\xi_t^N&\le -2\rho\hat\xi_t^N\md t \\
    &\qquad +2\left\langle\Delta\vartheta_t,D_\vartheta
    \left[
    \Psi_t(\bar\vartheta_t,G_{\bar\nu_t^{[N]}}(\bar\vartheta_t))
    -
    \Psi_t(\bar\vartheta_t,G_{\bar\nu_t}(\bar\vartheta_t))\right]\right\rangle \md t + \md M^N_t
\end{align*}
almost surely. By Cauchy--Schwarz and Assumption~\Cref{ass:psi_F_lipschitz},
\begin{align*}
&\left|
\left\langle
    \Delta\vartheta_t,
    D_\vartheta
    \left[
    \Psi_t(\bar\vartheta_t,G_{\bar\nu_t^{[N]}}(\bar\vartheta_t))
    -
    \Psi_t(\bar\vartheta_t,G_{\bar\nu_t}(\bar\vartheta_t))
    \right]
\right\rangle
\right|\\
&\qquad\le
    \|\Delta\vartheta_t\|_{D_{\vartheta}}
    \left\|
    \Psi_t(\bar\vartheta_t,G_{\bar\nu_t^{[N]}}(\bar\vartheta_t))
    -
    \Psi_t(\bar\vartheta_t,G_{\bar\nu_t}(\bar\vartheta_t))
    \right\|_{D_{\vartheta}}
    \\
    &\qquad\le
    \sqrt{\lambda_\vartheta^+}\sqrt{\widehat\xi_t^N}
    \left\|
    \Psi_t(\bar\vartheta_t,G_{\bar\nu_t^{[N]}}(\bar\vartheta_t))
    -
    \Psi_t(\bar\vartheta_t,G_{\bar\nu_t}(\bar\vartheta_t))
    \right\|\\
&\qquad\le
    L_\Psi\sqrt{\lambda_\vartheta^+}\,
    \sqrt{\widehat\xi_t^N}
    \left\|
    G_{\bar\nu_t^{[N]}}(\bar\vartheta_t)
    -
    G_{\bar\nu_t}(\bar\vartheta_t)
    \right\|,
\end{align*}

% \begin{align*}
% \frac{\mathrm d}{\mathrm dt}\mathbb E[\widehat\xi_t^N]
% &\le
% -2\rho\,\mathbb E[\widehat\xi_t^N]\\
% &\quad+
% 2\mathbb E
% \left[
% \left\langle
%     \Delta\vartheta_t,
%     D_\vartheta
%     \left[
%     \Psi_t(\bar\vartheta_t,G_{\bar\nu_t^{[N]}}(\bar\vartheta_t))
%     -
%     \Psi_t(\bar\vartheta_t,G_{\bar\nu_t}(\bar\vartheta_t))
%     \right]
% \right\rangle
% \right].
% \end{align*}

% \begin{align*}
% &\left|
% \mathbb E
% \left[
% \left\langle
%     \Delta\vartheta_t,
%     D_\vartheta
%     \left[
%     \Psi_t(\bar\vartheta_t,G_{\bar\nu_t^{[N]}}(\bar\vartheta_t))
%     -
%     \Psi_t(\bar\vartheta_t,G_{\bar\nu_t}(\bar\vartheta_t))
%     \right]
% \right\rangle
% \right]
% \right|\\
% &\qquad\le
%     \mathbb E[\|\Delta\vartheta_t\|^2_{D_{\vartheta}}]^{1/2}
%     \mathbb E
%     \left[
%     \left\|
%     \Psi_t(\bar\vartheta_t,G_{\bar\nu_t^{[N]}}(\bar\vartheta_t))
%     -
%     \Psi_t(\bar\vartheta_t,G_{\bar\nu_t}(\bar\vartheta_t))
%     \right\|_{D_{\vartheta}}^2
%     \right]^{1/2},\\
%     &\qquad\le
%     \sqrt{\lambda_\vartheta^+}\mathbb E\left[\widehat\xi_t^N\right]^{1/2}
%     \mathbb E
%     \left[
%     \left\|
%     \Psi_t(\bar\vartheta_t,G_{\bar\nu_t^{[N]}}(\bar\vartheta_t))
%     -
%     \Psi_t(\bar\vartheta_t,G_{\bar\nu_t}(\bar\vartheta_t))
%     \right\|^2
%     \right]^{1/2},\\
% &\qquad\le
%     L_\Psi\sqrt{\lambda_\vartheta^+}\,
%     \mathbb E\left[\widehat\xi_t^N\right]^{1/2}
%     \mathbb E
%     \left[
%     \left\|
%     G_{\bar\nu_t^{[N]}}(\bar\vartheta_t)
%     -
%     G_{\bar\nu_t}(\bar\vartheta_t)
%     \right\|^2
%     \right]^{1/2},
% \end{align*}
where \(\lambda_\vartheta^+\) is the largest eigenvalue of \(D_\vartheta\).
Since \(\bar\vartheta_0\) is deterministic, \(\bar\vartheta_t\) is
deterministic. Moreover, because the sampler initial states are i.i.d.\ and
independent of the mutually independent Brownian motions and each copy shares a deterministic optimiser path, the random vectors
\[
    Y_i
    :=
    F(\Pi_\theta(\bar\vartheta_t),\Pi_x(\bar Z_t^i))
    -
    G_{\bar\nu_t}(\bar\vartheta_t),
    \qquad i=1,\ldots,N,
\]
are independent and centred. Therefore, by
Assumption~\ref{ass:mf_bounded_second_moment},
\begin{align*}
    \mathbb E
    \left[
    \left\|
    G_{\bar\nu_t^{[N]}}(\bar\vartheta_t)
    -
    G_{\bar\nu_t}(\bar\vartheta_t)
    \right\|^2
    \right]
    &=
    \mathbb E\left\|\frac1N\sum_{i=1}^NY_i\right\|^2=
    \frac1{N^2}\sum_{i=1}^N\mathbb E\|Y_i\|^2
    \le
    \frac{V_F}{N}.
\end{align*}
Set
\[
    c:=L_\Psi\sqrt{\lambda_\vartheta^+},
    \qquad
    g_t:=\left\|
    G_{\bar\nu_t^{[N]}}(\bar\vartheta_t)
    -
    G_{\bar\nu_t}(\bar\vartheta_t)
    \right\|,
\]
so that \(\mathbb E[g_t^2]\le V_F/N\) for every \(t\ge0\) by the last display. Combining the almost-sure inequality for \(\md\widehat\xi_t^N\) with the Cauchy--Schwarz bound above gives, almost surely,
\[
    \md\widehat\xi_t^N
    \le
    \left(-2\rho\,\widehat\xi_t^N+2c\sqrt{\widehat\xi_t^N}\,g_t\right)\md t
    +\md M_t^N .
\]
Since \(t\mapsto e^{2\rho t}\) is deterministic and of finite variation, It\"o's product rule and \(\widehat\xi_0^N=0\) (synchronous initialisation) give, almost surely for all \(t\ge0\),
\[
    e^{2\rho t}\widehat\xi_t^N
    \le
    2c\int_0^t e^{2\rho s}\sqrt{\widehat\xi_s^N}\,g_s\,\md s
    +
    \int_0^t e^{2\rho s}\,\md M_s^N .
\]
Since no moment bounds for the particle system are assumed, the local martingale \(M^N\) need not be a true martingale, and we use the standard localisation argument for continuous local martingales, followed by Fatou's lemma to remove the stopping times; see for example \citet{stochastic_gronwal_lemma}. For \(n\ge1\), define the stopping times
\[
    \tau_n
    :=
    \inf\left\{s\ge0:\ \widehat\xi_s^N\ge n\ \text{or}\ \langle M^N\rangle_s\ge n\right\},
    \qquad
    \inf\emptyset:=\infty .
\]
Since \(\widehat\xi^N\) and \(\langle M^N\rangle\) are continuous, adapted and finite-valued, and \(\widehat\xi_0^N=0\), we have \(\tau_n\uparrow\infty\) almost surely as \(n\to\infty\), and \(\widehat\xi^N_{s\wedge\tau_n}\le n\) for all \(s\ge0\). Moreover, for fixed \(t\ge0\), the stopped integral \(\int_0^{t\wedge\tau_n}e^{2\rho s}\,\md M_s^N\) has quadratic variation at most \(e^{4\rho t}n\); hence it is a square-integrable martingale with zero mean. Evaluating the previous inequality at \(t\wedge\tau_n\), taking expectations and using Tonelli's theorem (the integrand is non-negative) gives
\[
    \phi_n(t)
    :=
    \mathbb E\left[e^{2\rho(t\wedge\tau_n)}\,\widehat\xi^N_{t\wedge\tau_n}\right]
    \le
    2c\int_0^t
    \mathbb E\left[\mathbf 1_{\{s<\tau_n\}}\,e^{2\rho s}\sqrt{\widehat\xi_s^N}\,g_s\right]\md s .
\]
Since \(\mathbf 1_{\{s<\tau_n\}}e^{2\rho s}\widehat\xi_s^N\le e^{2\rho(s\wedge\tau_n)}\widehat\xi^N_{s\wedge\tau_n}\), the Cauchy--Schwarz inequality and \(\mathbb E[g_s^2]\le V_F/N\) give, for every \(s\ge0\),
\[
    \mathbb E\left[\mathbf 1_{\{s<\tau_n\}}\,e^{2\rho s}\sqrt{\widehat\xi_s^N}\,g_s\right]
    \le
    \left(\mathbb E\left[\mathbf 1_{\{s<\tau_n\}}\,e^{2\rho s}\,\widehat\xi_s^N\right]\right)^{1/2}
    \left(\mathbb E\left[e^{2\rho s}g_s^2\right]\right)^{1/2}
    \le
    e^{\rho s}\sqrt{\frac{V_F}{N}}\,\sqrt{\phi_n(s)} .
\]
Hence \(\phi_n\) is continuous (by dominated convergence), satisfies \(0\le\phi_n(r)\le e^{2\rho r}n\), and, with \(a(s):=2c\sqrt{V_F/N}\,e^{\rho s}\),
\[
    \phi_n(r)\le\int_0^r a(s)\sqrt{\phi_n(s)}\,\md s
    \qquad\text{for all } r\ge0 .
\]
We now use the following square-root form of Gr\"onwall's lemma, a special case of the Bihari--LaSalle inequality: if \(\phi:[0,t]\to[0,\infty)\) is bounded and measurable, \(a\ge0\) is integrable on \([0,t]\), and \(\phi(r)\le\int_0^r a(s)\sqrt{\phi(s)}\,\md s\) for all \(r\in[0,t]\), then \(\sqrt{\phi(t)}\le\frac12\int_0^t a(s)\,\md s\). Indeed, \(\Phi(r):=\int_0^r a(s)\sqrt{\phi(s)}\,\md s\) is absolutely continuous and \(\phi\le\Phi\); for every \(\delta>0\), the function \(r\mapsto\sqrt{\Phi(r)+\delta}\) is absolutely continuous with
\[
    \frac{\md}{\md r}\sqrt{\Phi(r)+\delta}
    =
    \frac{a(r)\sqrt{\phi(r)}}{2\sqrt{\Phi(r)+\delta}}
    \le
    \frac{a(r)}{2}
    \qquad\text{for almost every } r\in[0,t],
\]
so that \(\sqrt{\phi(t)}\le\sqrt{\Phi(t)+\delta}\le\sqrt\delta+\frac12\int_0^t a(s)\,\md s\), and we let \(\delta\downarrow0\). Applying this lemma to \(\phi_n\) yields
\[
    \sqrt{\phi_n(t)}
    \le
    c\sqrt{\frac{V_F}{N}}\int_0^t e^{\rho s}\,\md s
    =
    \frac{c}{\rho}\sqrt{\frac{V_F}{N}}\left(e^{\rho t}-1\right).
\]
Since \(\tau_n\uparrow\infty\) almost surely, \(e^{2\rho(t\wedge\tau_n)}\widehat\xi^N_{t\wedge\tau_n}\to e^{2\rho t}\widehat\xi_t^N\) almost surely, and Fatou's lemma gives
\[
    e^{2\rho t}\,\mathbb E\left[\widehat\xi_t^N\right]
    \le
    \liminf_{n\to\infty}\phi_n(t)
    \le
    \frac{c^2}{\rho^2}\,\frac{V_F}{N}\left(e^{\rho t}-1\right)^2 .
\]
Recalling \(c=L_\Psi\sqrt{\lambda_\vartheta^+}\), we arrive at
\[
    \left(\mathbb E\left[\widehat\xi^N_t\right]\right)^{1/2}
    \le
    \frac{L_\Psi}{\rho}\sqrt{\frac{\lambda^+_\vartheta V_F}{N}}\left(1-e^{-\rho t}\right).
\]
%
% Therefore
% \[
%     \frac{\mathrm d}{\mathrm dt}
%     \mathbb E[\widehat\xi_t^N]
%     \le
%     -2\rho\,\mathbb E[\widehat\xi_t^N]
%     +
%     2L_\Psi
%     \sqrt{\frac{\lambda_\vartheta^+V_F}{N}}
%     \mathbb E[\widehat\xi_t^N]^{1/2}.
% \]
% Writing \(u_t:=\mathbb E[\widehat\xi_t^N]^{1/2}\), this implies
% \[
%     \frac{\mathrm d}{\mathrm dt}u_t
%     \le
%     -\rho u_t
%     +
%     L_\Psi
%     \sqrt{\frac{\lambda_\vartheta^+V_F}{N}}.
% \]
% The synchronous initialisation gives \(u_0=0\), and hence
% \[
%     u_t
%     \le
%     \frac{L_\Psi}{\rho}
%     \sqrt{\frac{\lambda_\vartheta^+V_F}{N}}
%     (1-e^{-\rho t}).
% \]
Finally,
\[
    W_2^2(\nu_t^N,\bar\nu_t^{[N]})
    \le
    \frac1N\sum_{i=1}^N\|Z_t^i-\bar Z_t^i\|^2,
\]
and
\[
    \xi_t^N
    \le
    \frac{1}{\min(\lambda_\vartheta^-,{\lambda_z^-})}
    \widehat\xi_t^N.
\]
Thus
\begin{align*}
    \Ltwonorm[\big]{\mathbf d\big((\vartheta_t^N,\nu_t^N),(\bar\vartheta_t,\bar\nu_t^{[N]})\big)}
    &\mathrel{\le}
    \left(\mathbb E\left[\xi_t^N\right]\right)^{1/2}
    \le
    \frac{L_\Psi}{\rho}
    \sqrt{
        \frac{\lambda_\vartheta^+V_F}
             {N\min(\lambda_\vartheta^-,{\lambda_z^-})}
    }
    (1-e^{-\rho t})\\
    &\le
    \frac{L_\Psi}{\rho}
    \sqrt{\frac{\kappa V_F}{N}}
    (1-e^{-\rho t}),
\end{align*}
which is the desired estimate.

\subsection{Proof of Theorem \ref{thm:parameter_error_full}}\label{app:cor_parameter_error_full}
Construct the \(N\)-copy mean-field
system and the interacting particle system with the coupling from
Theorem~\ref{thm:finite_particle_error}: set
\(\bar\vartheta_0=\vartheta_0^N\), let the mean-field sampler copies have
common initial law \(\nu_0\), and use the same initial sampler states and
Brownian motions in each particle--mean-field pair. Since every
mean-field sampler copy has law \(\bar\nu_t\), the common optimiser path
\(\bar\vartheta_t\) is the optimiser path of the one-copy mean-field system
started from \((\vartheta_0^N,\nu_0)\).

Write \(\bar\theta_t:=\Pi_\theta(\bar\vartheta_t)\), which is deterministic. Minkowski's inequality in \(\mathrm{L}^2(\mathbb P)\)
and the fact that \(\Pi_\theta\) is a Euclidean coordinate projection give
\[
\begin{aligned}
\Ltwonorm{\theta_t^N-\theta_\star}
&\leq
\Ltwonorm{\theta_t^N-\bar\theta_t}
+
\|\bar\theta_t-\theta_\star\|\\
&\leq
\Ltwonorm[\big]{
\mathbf d\big(
(\vartheta_t^N,\nu_t^N),
(\bar\vartheta_t,\bar\nu_t^{[N]})
\big)} +
\mathbf d\left(
(\bar\vartheta_t,\bar\nu_t),
(\vartheta_\star,\nu_{\theta_\star})
\right).
\end{aligned}
\]
By Theorem~\ref{thm:finite_particle_error}, the first term is bounded by
\[
\frac{L_\Psi}{\rho}
\sqrt{\frac{\kappa V_F}{N}}
(1-e^{-\rho t}).
\]
By Theorem~\ref{thm:invariant_measure},
\((\vartheta_\star,\nu_{\theta_\star})\) is a stationary mean-field
solution. Therefore
Theorem~\ref{thm:twisted_joint_strong_monotonicity} gives
\[
\begin{aligned}
&\mathbf d\left(
(\bar\vartheta_t,\bar\nu_t),
(\vartheta_\star,\nu_{\theta_\star})
\right)\leq
\sqrt\kappa\,e^{-\rho t}
\mathbf d\left(
(\bar\vartheta_0,\nu_0),
(\vartheta_\star,\nu_{\theta_\star})
\right)
=
\sqrt\kappa\,e^{-\rho t}\mathbf d_0,
\end{aligned}
\]
where \(\bar\vartheta_0=\vartheta_0^N\). Combining the last three displays
yields
\[
\Ltwonorm{\theta_t^N-\theta_\star}
\leq
\sqrt\kappa\,e^{-\rho t}\mathbf d_0
+
\frac{L_\Psi}{\rho}
\sqrt{\frac{\kappa V_F}{N}}
(1-e^{-\rho t}),
\]
which is the claimed parameter-error estimate.